\chapter{Theoretical Synthesis}
\label{ch:essay-synthesis}

Representation Invariance, Admissibility Preservation, Monotonic Relaxation, and Non-Identifiability are not four independent results requiring separate theoretical foundations. They are four consequences of a single admissibility relation examined at different points in its operation: Theorem~1.1 at a single state, Theorem~2.1 along a trajectory, Theorem~3.1 under restorative dynamics, and Theorem~4.2 under lossy projection. What varies is not the fundamental object of study but the operation being performed on it. This unity becomes consequential when one recognizes that the same structural mechanism can manifest as apparently different phenomena depending on where in the system it is observed.

Consider the reachable-set collapse that defines degeneracy in Theorem~2.1. A trajectory becomes incoherent when $\reach{S_i}$ collapses to the empty set or to a singleton already determined by prior history, because at that point the admissibility relation ceases to discriminate among continuations in any non-vacuous way. This is a forward-looking failure: the system has reached a state from which it cannot continue coherently. The same collapse, when examined retrospectively rather than prospectively, produces the observational collisions addressed by Theorem~4.2. If two distinct historical trajectories $H_1$ and $H_2$ both lead to states sharing the same reachable-set structure, then any observation map that records only what remains admissible from those states cannot distinguish which trajectory actually occurred. The forward-looking degeneracy and the backward-looking non-identifiability are therefore not separate failures but the same reachable-set structure examined from opposite temporal directions.

A concrete illustration clarifies this relationship. Suppose a discourse trajectory in which a speaker's utterances progressively narrow the space of admissible next contributions until only a single continuation remains licensed by what has been said so far. At that point the trajectory has degenerated in the sense of Theorem~2.1, because $\reach{S_t}$ is a singleton already fixed by $S_0, \dots, S_{t-1}$, and the admissibility relation is no longer performing discriminative work. Now consider an observer attempting to reconstruct how the discourse arrived at state $S_t$ from evidence consisting solely of which continuations remain admissible at that point. If two different discourse histories, one that narrowed possibilities through explicit assertion and another that narrowed them through presupposition and implicature, both produce the same singleton reachable set, then the observation map $\obsmap(H) = \reach{S_t}$ produces a collision: $\obsmap(H_1) = \obsmap(H_2)$ despite $H_1 \neq H_2$. The diagnostic corollary, Corollary~4.1, then applies: since the two histories share the same reachable set by construction, no observation map factoring through reachable-set structure can distinguish them, and the reconstruction failure is a structural limit rather than a correctable observational error.

The joint treatment thereby reveals that what appears as incoherence when examined forward through a trajectory and what appears as non-identifiability when examined backward through a reconstruction are related manifestations of reachable-set structure. The distinction is not merely that the same mathematical apparatus applies to both cases; it is that the collapse of discriminative structure causes both phenomena simultaneously. A system that has degenerated into a state from which only one continuation is admissible has also, by that very fact, entered a state from which its past becomes harder to recover, because the richness of the reachable set is precisely what would have allowed an observer to infer how that state was reached. Conversely, a pair of histories that produce observationally indistinguishable outcomes do so because they lead to states with identical reachable structure, and identical reachable structure is exactly the condition under which the forward trajectory from either state would be indistinguishable.

This relationship extends to the repair dynamics addressed by Theorem~3.1. Repair, defined as transitions that increase reachable-set volume, moves the system away from degeneracy by expanding $\mu(\reach{S_t})$. Since degeneracy is the condition under which both coherence fails and historical non-identifiability becomes more likely, repair trajectories that successfully increase reachable-set volume thereby simultaneously restore coherence in the forward direction and create observational distinctions that make historical reconstruction more feasible in the backward direction. The scalar energy functional $\closure{S_t} = -\log \mu(\reach{S_t})$ therefore measures not only inadmissibility in an abstract sense but also, concretely, the degree to which the system has approached the dual failures of forward incoherence and backward non-identifiability. Monotonic decrease of $\closure{S_t}$ along a repair trajectory signals movement away from both failure modes at once.

The cross-modal application developed in this essay benefits from this joint treatment because projection across modalities introduces opportunities for reachable-set structure to be preserved, altered, or destroyed in ways that affect representation, coherence, and recoverability simultaneously. If projecting semantic structure into speech preserves certain reachable-set relations while projecting the same structure into gesture preserves different relations, then Theorem~1.1 establishes that the two projections are representationally equivalent only to the extent that the preserved relations coincide. If either projection produces states with collapsed reachable sets, then Theorem~2.1 predicts that continuations in that modality will degenerate. If both projections lead to states sharing the same reachable structure despite originating from different semantic intentions, then Theorem~4.2 predicts that reconstructing the original intention from the projected signal will encounter non-identifiability. The four theorems thereby constrain the space of possible cross-modal behaviors jointly rather than imposing four independent constraints that must be checked separately.

Treating the four theorems as aspects of a single admissibility relation does not reduce them to triviality; rather, it allows their interactions to be studied systematically. A mechanism that generates behavior, whether control-based as discussed in Chapter~\ref{ch:essay-pct} or otherwise, operates within a space of trajectories that must satisfy the coherence condition of Theorem~2.1 if they are to avoid degeneracy. The observable projection of those trajectories is subject to the non-identifiability condition of Theorem~4.2 whenever the projection discards distinctions beyond those encoded in reachable-set structure. The representational equivalences licensed by Theorem~1.1 determine which different physical realizations of behavior count as semantically interchangeable. And the repair dynamics of Theorem~3.1 provide one account of how systems might navigate toward regions of state space where reachable sets remain nontrivial, thereby avoiding the dual failures of incoherence and non-identifiability. These are not four separate constraints applied in sequence but four views of the same structural condition examined from different angles, and recognizing them as such clarifies both what the admissibility framework establishes and what it leaves open for empirical investigation.
